\documentclass{article}

\usepackage{amsmath}
\usepackage{amssymb}
\usepackage{bussproofs}
\usepackage{fancyhdr}
\usepackage{mathpartir}

\pagestyle{fancy}
\fancyhf{} % Clear all header and footer fields

% Header: left, center, right
\fancyhead[L]{Programming Languages 2}
\fancyhead[C]{Homework 02}
\fancyhead[R]{Simon Sutoris 7542170}

% Horizontal line below the header
%\renewcommand{\headrulewidth}{0.4pt}

% Page number centered in the footer
\fancyfoot[C]{\thepage}

\begin{document}

\section*{Evaluation Contexts}

The upper rules are computation rules, the lower congruence rules.

Computation rules:

\[
\begin{array}{rlr}
v := & n & Values \\
v_1 \oplus v_2 \rightarrow & & (v_1 \oplus v_2) \\
\texttt{ifz}(0; t_1, t_2) \rightarrow & & t_1 \\
\texttt{ifz}(v; t_2, t_2) \rightarrow & & t_2\ \text{if}\ n \ne 0
\end{array}
\]

Congruence rules:

\[
\begin{array}{ccc}
\inferrule{t_1 \rightarrow t'_1}{t_1 \oplus t_2 \rightarrow t'_1 \oplus t_2} &
\inferrule{t_2 \rightarrow t'_2}{v \oplus t_2 \rightarrow v \oplus t'_2} &
\inferrule{t_1 \rightarrow t'_1}{\texttt{ifz}(t_1; t_2, t_3) \rightarrow \texttt{ifz}(t'_1; t_2, t_3)}
\end{array}
\]

Congruence rules re-formulated using \textit{evaluation contexts}:

\[
\begin{array}{ll}
E ::= & \\
& [\ ] \\
& E \oplus t \\
& v \oplus E \\
& \texttt{ifz}(E; t_2, t_3)
\end{array}
\]

\section*{Evaluation}

\subsection*{Derivation Trees}

%\[
%\inferrule{(3 * 4) \rightarrow 12}{\texttt{ifz}((3 * 4) - 5; (7 * 7) + 15, 4 + 4) 7 \rightarrow \texttt{ifz}(12 - 5; (7 * 7) + 15, 4 + 4)}
%\]

\begin{prooftree}
  \AxiomC{3 * 4 \to 12}
  \UnaryInfC{(3 * 4) - 5 \to 12 - 5}
  \UnaryInfC{\texttt{ifz}((3 * 4) - 5; (7 * 7) + 15, 4 + 4) 7 \to \texttt{ifz}(12 - 5; (7 * 7) + 15, 4 + 4)}
\end{prooftree}

\subsection*{Normal Forms}

A term which is in normal form and is a value is quite simple: $t_1 = 42$.

However, a term which is in normal form and not a value does not exist. The grammer of the language only allows for three terms: values, concatenation between terms using binary operations, or \texttt{ifz} which consumes 3, possibly different terms.

To not be in normal form, a term has to be found which cannot be evaluated. A value would fit this requirement, but the question also asks for a term that is not a value. Therefore, a concatenation between binary operations or usage of \texttt{ifz} has to be found.

Let's first look at the binary operations. Either $t_1$ is not a value, congruence rule 1 applies.

\subsection*{Induction on Derivation Trees}

\subsection*{(i) $t = v$}

\begin{align}
\texttt{size}(t) = \texttt{size}(v) = 1 = \texttt{size(t')}
\end{align}

\subsection*{(ii) $t = t_1 \oplus t_2, t_1 \to^* v_1$}

\begin{align}
\texttt{size}(t_1 \oplus t_2) 
&= \texttt{size}(t_1) + \texttt{size}(t_2) + 1 \label{eq:case-2-1}\\
&\ge \texttt{size}(v_1) + \texttt{size}(t_2) + 1 \label{eq:case-2-2}\\
&= \texttt{size}(t')\\
\Leftrightarrow \texttt{size}(t_1) &\ge \texttt{size}(v_1) \\
\texttt{size}(t_1) &\ge 1
\end{align}

Analog for $t = t_1 \oplus t_2, t_2 \to^* v_2$

\subsection*{(iii) $t = \texttt{ifz}(t_1; t_2, t_3), t_1 \to^* v_1$}

\begin{align}
\texttt{size}(\texttt{ifz}(t_1; t_2, t_3))
&= \texttt{size}(t_1) + \texttt{size}(t_2) + \texttt{size}(t_3) + 1 \\
&\ge \texttt{size}(v_1) + \texttt{size}(t_2) + \texttt{size}(t_3) + 1\\
&= \texttt{size}(t')\\
\Leftrightarrow \texttt{size}(t_1) &\ge \texttt{size}(v_1) \\
\texttt{size}(t_1) &\ge 1
\end{align}

\section*{Types and Programming Languages}

Is OCaml mandatory or can we use it at home to veryfy the contents of the book? Also, are the questions mandatory? Chapter 4 was rather short compared to the beginnings of the book till Chapter 3, and I feel like I have to force myself to come up with a question?

\end{document}
